% gz-p8-01-1: 事件运算 Venn 图——A、B、A∪B、A∩B
\documentclass[tikz,border=4pt]{standalone}
\usepackage{ctex}
\usepackage{amsmath}
\usetikzlibrary{calc}

\begin{document}
\begin{tikzpicture}[scale=1.0, thick]
  % 左：A ∪ B（涂色整体）
  \begin{scope}[shift={(0,0)}]
    % 矩形 Ω
    \draw[black] (-2, -2) rectangle (3, 2.4);
    \node[above right, font=\small] at (-2, 2.4) {$\Omega$};

    % 涂色：先涂 A
    \fill[red!30] (0,0) circle (1.4);
    % 再涂 B
    \fill[red!30] (1.5, 0) circle (1.4);

    % 边界
    \draw[black, thick] (0,0) circle (1.4);
    \draw[black, thick] (1.5, 0) circle (1.4);

    % 标签
    \node[font=\small] at (-0.7, 1.1) {$A$};
    \node[font=\small] at (2.2, 1.1) {$B$};

    \node[below, font=\small] at (0.75, -2.2) {$A\cup B$（涂色）};
  \end{scope}

  % 右：A ∩ B（涂色交集）
  \begin{scope}[shift={(7,0)}]
    % 矩形 Ω
    \draw[black] (-2, -2) rectangle (3, 2.4);
    \node[above right, font=\small] at (-2, 2.4) {$\Omega$};

    % 先画圆
    \draw[black, thick] (0,0) circle (1.4);
    \draw[black, thick] (1.5, 0) circle (1.4);

    % 涂色：交集
    \begin{scope}
      \clip (0,0) circle (1.4);
      \fill[blue!40] (1.5, 0) circle (1.4);
    \end{scope}

    % 重画边界
    \draw[black, thick] (0,0) circle (1.4);
    \draw[black, thick] (1.5, 0) circle (1.4);

    % 标签
    \node[font=\small] at (-0.7, 1.1) {$A$};
    \node[font=\small] at (2.2, 1.1) {$B$};
    \node[font=\footnotesize, blue] at (0.75, 0) {$A\cap B$};

    \node[below, font=\small] at (0.75, -2.2) {$A\cap B$（涂色）};
  \end{scope}
\end{tikzpicture}
\end{document}
